\documentclass[11pt,a4paper]{article}

\usepackage[margin=2.2cm]{geometry}
\usepackage[T1]{fontenc}
\usepackage[utf8]{inputenc}
\usepackage{lmodern}
\usepackage{microtype}
\usepackage[hidelinks]{hyperref}
\usepackage{titlesec}
\usepackage{enumitem}

\titleformat{\section}{\normalsize\bfseries}{\thesection.}{0.5em}{}
\titlespacing*{\section}{0pt}{1.0ex}{0.5ex}
\setlist{itemsep=1pt,topsep=2pt,leftmargin=1.2em}
\setlength{\parskip}{3pt}

\title{\vspace{-1.5cm}\large\textbf{Reading Assignment: Tracking Aviation's Footprint from the Sky}\\[2pt]
\normalsize Critical review of \emph{``Global Civil Aviation Emissions Estimates for 2017--2020 Using ADS-B Data''}\\
(Quadros, Snellen, Sun \& Dedoussi, \emph{Journal of Aircraft}, 2022)}
\author{Pedro Gonz\'alez N\'u\~nez \and Diego Colin Reyes \and Nacho Castilla Llorca}
\date{}

\begin{document}
\maketitle
\vspace{-1.2cm}

\section{Problem}
Aviation is responsible for roughly 2--3\% of global anthropogenic CO$_2$ emissions, and its total
climate effect is larger once non-CO$_2$ effects such as NO$_x$ and contrails are included~\cite{lee2021}.
Despite this, our knowledge of \emph{where} and \emph{when} aircraft emit is surprisingly limited.
Official figures (ICAO, IEA) are mostly \emph{top-down}: they divide jet-fuel sales between countries,
come out with a delay of a year or more, and say nothing about the altitude or location of emissions,
which atmospheric models need. Earlier \emph{bottom-up} inventories, built from flight schedules or
radar data, are either proprietary, regionally limited, or not updated every year. This became a clear
problem in 2020: when COVID-19 grounded much of the world's fleet, there was no timely, open and
spatially resolved estimate of how much aviation emissions actually fell. The question the authors ask
is whether the huge volume of aircraft position data now broadcast via ADS-B (Automatic Dependent
Surveillance--Broadcast) can be turned into a reliable, up-to-date global emissions inventory~\cite{quadros2022}.

\section{Approach}
The authors build a bottom-up pipeline that turns raw flight observations into gridded, speciated emissions:
\begin{itemize}
  \item \textbf{Activity data.} Flights observed by crowdsourced networks of ground ADS-B receivers are
  aggregated into flight lists of (aircraft type, origin, destination, count) for 2017--2020.
  \item \textbf{Trajectory model.} Each origin--destination pair is flown along the great circle,
  extended by a constant lateral-inefficiency term (40.5\,NM + 3.87\% of the distance), with cruise
  levels set from observed data and monthly-averaged winds.
  \item \textbf{Fuel burn.} Climb, cruise and descent are simulated with EUROCONTROL's BADA~3.15
  aircraft performance model; the landing--takeoff cycle uses a time-in-mode approach.
  \item \textbf{Emissions.} CO$_2$, H$_2$O and SO$_x$ scale with fuel; NO$_x$, CO and HC use ICAO Engine
  Emissions Databank indices corrected for altitude (Boeing Fuel Flow Method~2); non-volatile particulate
  matter (nvPM) is estimated from engine smoke numbers.
  \item \textbf{Openness.} The model is released as open-source software (openAVEM) and the gridded
  monthly results as a CC~BY~4.0 dataset.
\end{itemize}
\textbf{Key results.} For 2019 the model estimates about 297\,Tg of fuel, 937\,Tg of CO$_2$ and 4.62\,Tg of
NO$_x$. COVID-19 cut 2020 fuel burn by about 48\% relative to expected levels, avoiding roughly 463\,Tg of
CO$_2$ and 2.29\,Tg of NO$_x$. The authors also track changes in fleet-average emission indices over time
and conclude that ADS-B is a viable source for \emph{timely and transparent} global emissions monitoring.

\section{Critical Reflection}
\textbf{Where we agree.} The central idea is convincing and, in our view, important. The paper shows that a
data stream originally built for air-traffic safety can be reused as an environmental sensor. Producing a
global inventory within months instead of years, with full 3-D resolution, is a real step forward: it made
it possible to quantify the COVID-19 ``natural experiment'' almost in real time. Its estimates also hold
up well against independent work: the later GAIA inventory, which uses full satellite and terrestrial ADS-B
trajectories, finds 2019 fuel burn only 4.7\% lower and an almost identical 2020 drop
($-48.2\%$ vs.\ $-47.1\%$)~\cite{teoh2024}. Releasing the code and gridded data is good scientific practice
and makes the work useful to climate modellers.

\textbf{Where we are more critical.}
\begin{enumerate}
  \item \emph{``Transparent'' is only partly true.} The pipeline depends on BADA~3, which must be licensed from
  EUROCONTROL, and on flight lists derived from a commercial tracking platform~\cite{teoh2024}. An outside
  group cannot fully reproduce the inventory from public inputs. A truly open version would need open
  performance models (e.g.\ OpenAP) and open data (e.g.\ the OpenSky Network~\cite{schafer2014}).
  \item \emph{Coverage bias.} Ground receivers are dense in Europe and North America but sparse over oceans,
  Africa and parts of Asia. Missing flights simply do not appear, which helps explain why the result (297\,Tg)
  stays below the IEA's top-down 320\,Tg~\cite{teoh2024}. This bias is not neutral: it under-represents
  exactly the regions where air traffic is growing fastest.
  \item \emph{The data is richer than the model.} ADS-B provides real trajectories, yet the method reduces
  them to origin--destination counts and re-flies them as idealised great circles with monthly winds.
  This throws away actual routing, holding patterns, airspace closures and conflict-zone detours;
  GAIA shows that a fixed 4--6\% inefficiency underestimates distance flown in regions like Europe~\cite{teoh2024}.
  \item \emph{Emissions are not climate impact.} The paper reports emitted mass, but contrail cirrus may cause
  a larger warming effect than aviation's CO$_2$ itself~\cite{lee2021}. Location-aware ADS-B data is exactly what is
  needed to assess contrails, so stopping at emissions leaves the most policy-relevant question open.
  nvPM, which drives contrail formation, is also likely underestimated~\cite{teoh2024}.
\end{enumerate}

\textbf{Overall position.} We agree with the paper's main conclusion: ADS-B data \emph{can} support timely,
global, spatially resolved emissions monitoring, and the COVID-19 results are robust. We disagree with how
strongly it presents the approach as transparent, and we think its simplified trajectories under-use the very
data that makes it novel. Our own project, which ingests OpenSky state vectors into a Spark/MongoDB pipeline
to detect holding patterns and congestion, illustrates the next step: using real trajectories, not idealised
ones, to attribute emissions to specific inefficiencies that operational changes could actually remove.

\vspace{-0.3cm}
{\small
\begin{thebibliography}{9}\setlength{\itemsep}{0pt}
\bibitem{quadros2022} F.~D.~A.~Quadros, M.~Snellen, J.~Sun, I.~C.~Dedoussi, ``Global civil aviation emissions
estimates for 2017--2020 using ADS-B data,'' \emph{Journal of Aircraft}, 59(6):1394--1405, 2022.
doi:10.2514/1.C036763.
\bibitem{teoh2024} R.~Teoh, Z.~Engberg, M.~Shapiro, L.~Dray, M.~E.~J.~Stettler, ``The high-resolution Global
Aviation emissions Inventory based on ADS-B (GAIA) for 2019--2021,'' \emph{Atmos.\ Chem.\ Phys.}, 24:725--744, 2024.
\bibitem{lee2021} D.~S.~Lee et al., ``The contribution of global aviation to anthropogenic climate forcing for
2000 to 2018,'' \emph{Atmospheric Environment}, 244:117834, 2021.
\bibitem{schafer2014} M.~Sch\"afer, M.~Strohmeier, V.~Lenders, I.~Martinovic, M.~Wilhelm, ``Bringing up OpenSky:
A large-scale ADS-B sensor network for research,'' \emph{Proc.\ IPSN}, 2014.
\end{thebibliography}
}

\end{document}
